\chapter[El Cofre del Egoísmo]{El Cofre del Egoísmo\\ \large El corazón que no da, no tiene donde recibir.}

\elegantimage{09_frio_egoismo/web/assets/hero.png}

\lettrine[lines=3, nindent=0.5em, findent=0.2em]{N}{}egar la ayuda cuando se tiene la capacidad de darla es cerrar los canales por los que fluye la abundancia. El egoísta vive en la ilusión de posesión, sin entender que todo es energía en movimiento.

\section{El Frío de las Monedas}
\begin{elegantquote}
\textit{Un rico mercader se sentaba sobre montañas de monedas de oro, mientras cerraba la puerta a niños descalzos que tiritaban en la nieve. ``Lo mío es mío'', repetía, creyendo que su tesoro lo protegería de la muerte y del dolor...}

\textit{Pero el egoísmo congela el alma. Al encadenar su riqueza, se separó de la gran familia humana. Tiempo después, la rueda de la vida giró, y él se encontró a sí mismo siendo un vagabundo perdido en una calle sin nombre. Nadie lo miraba. Nadie le daba calor.}

\textit{Aquel mendigo que sufría el terrible frío del abandono, era él experimentando el mismo hielo que creó. Solo rompiéndose en llanto al descubrir el dolor del hambre, logrará descongelar su corazón... para comprender que el mayor tesoro no es retener, sino compartir el pan bajo el mismo techo del amor.}

\end{elegantquote}

\vspace{1em}
\begin{center}
    \textbf{\Large La escasez es el eco del corazón cerrado.}
\end{center}
\vspace{1em}

\section{Análisis Kármico y Traducción}
\begin{wrapfigure}{r}{0.5\textwidth}
  \vspace{-1.5em}
  \begin{center}
  \inlineelegantimage[0.45\textwidth]{09_frio_egoismo/web/assets/pasaje_original.jpg}
  \end{center}
  \vspace{-1.5em}
\end{wrapfigure}
El egoísmo y la tacañería aíslan al alma en una celda de frío espiritual, obligándola a renacer en condiciones de abandono absoluto para que finalmente aflore la sed de dar.

